% META: {"id":"1SPE-VARALEA-CO-022","chapitre":"1SPE-VARIABLES-ALEATOIRES","type_objet":"corrige","exercice_ref":"1SPE-VARALEA-EX-022","capacites_codes":[],"extension_codes":["X1"],"programme_alignment":"OPTIONAL_EXTENSION","extension_label":"Approfondissement — Vers la Terminale","sources_inspiration":[],"status":"approved","origin":"generated","status_history":["generated","audited","accepted","approved"],"release_acceptance":"RELEASE_OWNER_BATCH_ACCEPTANCE","acceptance_closure_digest":"sha256:9b3ccf9a81c5520fbb7e03b7b2d3e7bf2057a02834b6d908bf3be5f49f3b553f","acceptance_receipt_digest":"sha256:4fad86f0282e0fa4e0419cca1ad6379ae7af5a280c04a4a90880f90a1c044242"}
\begin{center}\textbf{Approfondissement — Vers la Terminale}\end{center}
% BEGIN-VERIFY
% from sympy import binomial, Rational
% p = Rational(1,2)
% P3 = binomial(5,3)*p**3*(1-p)**2
% assert P3 == Rational(5,16)
% P0 = p**5; assert P0 == Rational(1,32)
% P5 = p**5; assert P5 == Rational(1,32)
% END-VERIFY
\begin{corrige}{1SPE-VARALEA-EX-022}

\textbf{1.} Chaque lancer est une épreuve de Bernoulli ($p = \frac{1}{2}$), les $5$ lancers sont indépendants, $X$ compte le nombre de succès : $X \sim \mathcal{B}(5, \frac{1}{2})$.

\medskip
\textbf{2.} $P(X=3) = \binom{5}{3}\left(\frac{1}{2}\right)^3\left(\frac{1}{2}\right)^2 = 10 \times \frac{1}{32} = \frac{10}{32} = \frac{5}{16}$.

\medskip
\textbf{3.} $P(X=0) = \left(\frac{1}{2}\right)^5 = \frac{1}{32}$. \quad $P(X=5) = \left(\frac{1}{2}\right)^5 = \frac{1}{32}$.

\end{corrige}
